% GENERATED FILE. Edit canonical lessons, then run npm run generate:books.
% canonical-source-hash: fnv1a64:16d3ada72a044006
% canonical-lessons: AR-C476-aqrab, AR-C476-jarada, AR-C476-tuhlub, AR-C476-mikhlab, AR-C476-iftarasa

\chapter{Small creatures}
\label{ch:ar-small-creatures}
\hlchaptermodality{476}

\begin{chapteropening}
\textbf{By the end of this chapter:} \emph{I can say a scorpion, a locust, moss, a claw and to prey on.}

Complete the last lesson of chapter 476: I can say a scorpion, a locust, moss, a claw and to prey on.
\end{chapteropening}

\section[\texorpdfstring{\ar{عقرب} (ʿaqrab)}{عقرب (ʿaqrab)}]{\ar{عقرب} (ʿaqrab) — a scorpion}

\label{lesson:AR-C476-aqrab}



\begin{quote}
Before the new one: say the Arabic for a squirrel, then the Arabic for a bat.
\end{quote}

\subsection*{You'll want to know: \ar{عقرب}}
\textbf{\ar{عقرب}} — \emph{ʿaqrab} — \textquotedblleft{}a scorpion\textquotedblright{}.

\begin{quote}
\ar{وجدنا} \ar{عقربا} \ar{تحت} \ar{حجر}.
\end{quote}

\begin{quote}
\emph{wajadnā ʿaqraban taḥta ḥajar.}
\end{quote}

\begin{quote}
We found a scorpion under a stone.
\end{quote}

Plural: \ar{عقارب}, \emph{ʿaqārib}.

The hands of a clock are also \emph{ʿaqārib}.

\begin{cousinweb}
\textbf{\ar{عقرب}} (\emph{ʿaqrab}) goes back to an old Semitic word, with relatives in Hebrew, Syriac and Ugaritic (EtymDB).
\end{cousinweb}

\begin{grammarlens}[title={What you've built}]
One more word for getting around.
\end{grammarlens}

\subsection*{Your turn}
\begin{itemize}
\raggedright
  \item \emph{Say it:} \emph{ʿaqrab}
  \item \emph{Say it:} \emph{ʿaqrab}, once more
  \item \emph{Say it:} \emph{wajadnā ʿaqraban taḥta ḥajar.} — We found a scorpion under a stone.
  \item \emph{Recall:} say the Arabic for a squirrel, then the Arabic for a bat, then say \emph{ʿaqrab} again
\end{itemize}

\begin{tcolorbox}[breakable,colback=teal!4,colframe=teal!35!black,title={Before you move on}]
What does \ar{عقرب} mean? (\textquotedblleft{}A scorpion\textquotedblright{}.) Say it once more.
\end{tcolorbox}
\section[\texorpdfstring{\ar{جرادة} (jarāda)}{جرادة (jarāda)}]{\ar{جرادة} (jarāda) — a locust}

\label{lesson:AR-C476-jarada}



\begin{quote}
Before the new one: say the Arabic for a bat, then the Arabic for a scorpion.
\end{quote}

\subsection*{You'll want to know: \ar{جرادة}}
\textbf{\ar{جرادة}} — \emph{jarāda} — \textquotedblleft{}a locust\textquotedblright{}.

\begin{quote}
\ar{رأيت} \ar{جرادة} \ar{خضراء} \ar{على} \ar{الشجرة}.
\end{quote}

\begin{quote}
\emph{raʾaytu jarāda khaḍrāʾ ʿalā ash-shajara.}
\end{quote}

\begin{quote}
I saw a green locust on the tree.
\end{quote}

\emph{jarād} is locusts as a swarm; \emph{jarāda} is one.

\begin{cousinweb}
\textbf{\ar{جرادة}} (\emph{jarāda}) is built on the native Arabic root \emph{j-r-d}, \textquotedblleft{}to strip bare\textquotedblright{} (Lane). Locusts strip the land bare.
\end{cousinweb}

\begin{grammarlens}[title={What you've built}]
One more word for getting around.
\end{grammarlens}

\subsection*{Your turn}
\begin{itemize}
\raggedright
  \item \emph{Say it:} \emph{jarāda}
  \item \emph{Say it:} \emph{jarāda}, once more
  \item \emph{Say it:} \emph{raʾaytu jarāda khaḍrāʾ ʿalā ash-shajara.} — I saw a green locust on the tree.
  \item \emph{Recall:} say the Arabic for a bat, then the Arabic for a scorpion, then say \emph{jarāda} again
\end{itemize}

\begin{tcolorbox}[breakable,colback=teal!4,colframe=teal!35!black,title={Before you move on}]
What does \ar{جرادة} mean? (\textquotedblleft{}A locust\textquotedblright{}.) Say it once more.
\end{tcolorbox}
\section[\texorpdfstring{\ar{طحلب} (ṭuḥlub)}{طحلب (ṭuḥlub)}]{\ar{طحلب} (ṭuḥlub) — moss; algae}

\label{lesson:AR-C476-tuhlub}



\begin{quote}
Before the new one: say the Arabic for a scorpion, then the Arabic for a locust.
\end{quote}

\subsection*{You'll want to know: \ar{طحلب}}
\textbf{\ar{طحلب}} — \emph{ṭuḥlub} — \textquotedblleft{}moss; algae\textquotedblright{}.

\begin{quote}
\ar{كان} \ar{الطحلب} \ar{على} \ar{الصخور}.
\end{quote}

\begin{quote}
\emph{kāna aṭ-ṭuḥlub ʿalā aṣ-ṣukhūr.}
\end{quote}

\begin{quote}
There was moss on the rocks.
\end{quote}

\begin{cousinweb}
\textbf{\ar{طحلب}} (\emph{ṭuḥlub}) names moss and the green film on water; the origin is not settled.
\end{cousinweb}

\begin{grammarlens}[title={What you've built}]
One more word for getting around.
\end{grammarlens}

\subsection*{Your turn}
\begin{itemize}
\raggedright
  \item \emph{Say it:} \emph{ṭuḥlub}
  \item \emph{Say it:} \emph{ṭuḥlub}, once more
  \item \emph{Say it:} \emph{kāna aṭ-ṭuḥlub ʿalā aṣ-ṣukhūr.} — There was moss on the rocks.
  \item \emph{Recall:} say the Arabic for a scorpion, then the Arabic for a locust, then say \emph{ṭuḥlub} again
\end{itemize}

\begin{tcolorbox}[breakable,colback=teal!4,colframe=teal!35!black,title={Before you move on}]
What does \ar{طحلب} mean? (\textquotedblleft{}Moss; algae\textquotedblright{}.) Say it once more.
\end{tcolorbox}
\section[\texorpdfstring{\ar{مخلب} (mikhlab)}{مخلب (mikhlab)}]{\ar{مخلب} (mikhlab) — a claw; a talon}

\label{lesson:AR-C476-mikhlab}



\begin{quote}
Before the new one: say the Arabic for a locust, then the Arabic for moss.
\end{quote}

\subsection*{You'll want to know: \ar{مخلب}}
\textbf{\ar{مخلب}} — \emph{mikhlab} — \textquotedblleft{}a claw; a talon\textquotedblright{}.

\begin{quote}
\ar{أمسك} \ar{الصقر} \ar{السمكة} \ar{بمخالبه}.
\end{quote}

\begin{quote}
\emph{amsaka aṣ-ṣaqr as-samaka bi-makhālibihi.}
\end{quote}

\begin{quote}
The falcon held the fish in its talons.
\end{quote}

Plural: \ar{مخالب}, \emph{makhālib}.

\begin{cousinweb}
\textbf{\ar{مخلب}} (\emph{mikhlab}) is built on the native Arabic root \emph{kh-l-b}, \textquotedblleft{}to wound with the claws\textquotedblright{} (Lane). The tool pattern \emph{mifʿal} makes it the thing one claws with.
\end{cousinweb}

\begin{grammarlens}[title={What you've built}]
One more word for getting around.
\end{grammarlens}

\subsection*{Your turn}
\begin{itemize}
\raggedright
  \item \emph{Say it:} \emph{mikhlab}
  \item \emph{Say it:} \emph{mikhlab}, once more
  \item \emph{Say it:} \emph{amsaka aṣ-ṣaqr as-samaka bi-makhālibihi.} — The falcon held the fish in its talons.
  \item \emph{Recall:} say the Arabic for a locust, then the Arabic for moss, then say \emph{mikhlab} again
\end{itemize}

\begin{tcolorbox}[breakable,colback=teal!4,colframe=teal!35!black,title={Before you move on}]
What does \ar{مخلب} mean? (\textquotedblleft{}A claw; a talon\textquotedblright{}.) Say it once more.
\end{tcolorbox}
\section[\texorpdfstring{\ar{افترس} (iftarasa)}{افترس (iftarasa)}]{\ar{افترس} (iftarasa) — to prey on}

\label{lesson:AR-C476-iftarasa}



\begin{quote}
Before the new one: say the Arabic for moss, then the Arabic for a claw.
\end{quote}

\subsection*{You'll want to know: \ar{افترس}}
\textbf{\ar{افترس}} — \emph{iftarasa} — \textquotedblleft{}to prey on\textquotedblright{}.

\begin{quote}
\ar{يفترس} \ar{النسر} \ar{الأرانب}.
\end{quote}

\begin{quote}
\emph{yaftarisu an-nasr al-arānib.}
\end{quote}

\begin{quote}
The eagle preys on rabbits.
\end{quote}

Present: \ar{يفترس}, \emph{yaftaris}.

\begin{cousinweb}
\textbf{\ar{افترس}} (\emph{iftarasa}) is built on the native Arabic root \emph{f-r-s}, \textquotedblleft{}to break the neck of prey\textquotedblright{} (Lane). It is Form VIII; \emph{faras}, \textquotedblleft{}a horse\textquotedblright{}, is on the same root.
\end{cousinweb}

\begin{grammarlens}[title={What you've built}]
One more word for getting around.
\end{grammarlens}

\subsection*{Your turn}
\begin{itemize}
\raggedright
  \item \emph{Say it:} \emph{iftarasa}
  \item \emph{Say it:} \emph{iftarasa}, once more
  \item \emph{Say it:} \emph{yaftarisu an-nasr al-arānib.} — The eagle preys on rabbits.
  \item \emph{Recall:} say the Arabic for moss, then the Arabic for a claw, then say \emph{iftarasa} again
\end{itemize}

\begin{tcolorbox}[breakable,colback=teal!4,colframe=teal!35!black,title={Before you move on}]
What does \ar{افترس} mean? (\textquotedblleft{}To prey on\textquotedblright{}.) Say it once more.
\end{tcolorbox}
